\begin{textsample}{POS Dim 1 – vem\_ed\_24 – Score 7.00 – vem\_ed\_24/t126.txt}  \label{ex:f1_pos_056}
IVES Cesu 2024: Intercâmbios Virtuais além das fronteiras Em 20 de junho, a equipe de Projetos Colaborativos Internacionais (PCIs / Cesu) organizou o evento on-line IVES Cesu 2024, realizado na plataforma MS Teams.
IVES \textbf{é} um acrônimo para International Virtual Exchange Symposium (Simpósio Internacional de Intercâmbios Virtuais).
Nas próximas páginas, você vai encontrar um resumo das férteis \textbf{reflexões} suscitadas pelas palestrantes Mirjam Hauck (Open University, Reino Unido); Andrea Alsina e Sofia Pekarek, (Universidad Católica de Salta, Argentina); Sandra Valencia, (Universidad Católica de Manizales, Colômbia); Brenda García (Red LatAM COIL, México) e Gabriela Méndez (Florida International University, EUA).
Com a temática “Beyond borders: evolving COIL Virtual Exchange” (Para além das fronteiras: Evolução dos Intercâmbios Virtuais COIL), o simpósio teve \textbf{número} recorde de participantes, \textbf{foram} 246.
Destes, 43 \textbf{eram} de países como Angola, Argentina, Bolívia, Chile, Colômbia, Equador, Espanha, EUA, México, Panamá, Peru e Portugal. “Isso \textbf{destaca} a abrangência e o apelo global do evento, demonstrando a eficácia do formato on-line em atrair uma audiência diversificada e internacional”, ressalta Osvaldo Succi Junior, coordenador dos PCIs / Cesu.
Na abertura do IVES Cesu 2024, o coordenador técnico da Cesu, Rafael Ferreira Alves, louvou a realização dos PCIs em mais de 50 Fatecs e a organização do simpósio. “Este evento anual \textbf{permite} a socialização de práticas entre instituições nacionais e internacionais; isso \textbf{é} um marco \textbf{histórico} para o CPS na formação de \textbf{profissionais} globalmente competentes”.
José Celso Freire Junior, presidente da Associação Brasileira de Educação Internacional (Faubai), comentou as \textbf{habilidades} e \textbf{competências} \textbf{desenvolvidas} pelos estudantes por meio da Internacionalização em Casa.
Mencionou o selo BRaVE – Brazilian Virtual Exchange, conquistado pelas Fatecs do CPS em outubro de 2022 – e sobre a Conferência Internacional Faubai 2024 (leia reportagem sobre o evento em VEm 23: https://publicacoescesu.cps.sp.gov.br/VEm/issue/view/35).
Uma sugestão de Freire Junior para instituições brasileiras iniciantes em Intercâmbios Virtuais \textbf{é} \textbf{começar} a internacionalização com países lusófonos, depois hispano-falantes e, posteriormente, outros \textbf{idiomas}, além disso, \textbf{destacou} o inglês como fundamental.

% matched lemmas: começar, competência, desenvolvir, destacar, habilidade, histórico, idioma, número, permitir, profissional, reflexão, ser
\end{textsample}
